\chapter{Mathematical Foundations}
\label{chap:mathematical_foundations}

\begin{tldr}
This chapter covers the mathematical foundations essential for deep learning: linear algebra (matrix operations, eigenvalues, SVD), calculus (gradients, chain rule, Jacobians), probability (Bayes' theorem, distributions, MLE), and information theory (entropy, KL divergence, mutual information). These concepts underpin every architecture and training algorithm discussed in this guide.
\end{tldr}

\section{Linear Algebra for Deep Learning}
\label{sec:linear_algebra}

\subsection{Vectors and Matrix Operations}

\subsubsection{Key Operations}

\textbf{Dot Product.} For vectors $\vx, \vy \in \R^n$:
\begin{equation}
\vx \cdot \vy = \vx^T \vy = \sum_{i=1}^n x_i y_i = \|\vx\| \|\vy\| \cos\theta
\end{equation}
The geometric interpretation (cosine of the angle) is why cosine similarity works for embeddings.

\textbf{Applications in DL}:
\begin{itemize}
    \item Similarity computation (embedding dot product)
    \item Attention scores: $\vq^T \vk$
    \item Linear layer output: $\vh = \mW \vx + \vb$
\end{itemize}

\subsubsection{Matrix Multiplication}

For $\mA \in \R^{m \times n}$ and $\mB \in \R^{n \times p}$:
\begin{equation}
(\mA\mB)_{ij} = \sum_{k=1}^n A_{ik} B_{kj}
\end{equation}

\textbf{Computational complexity}: $O(mnp)$

\textbf{Key insight}: Matrix multiplication is the core operation in neural networks. Understanding its cost is essential for analyzing model efficiency.

\subsection{Eigenvalues and Eigenvectors}

\textbf{Eigenpair.} For matrix $\mA \in \R^{n \times n}$, scalar $\lambda$ and vector $\vv \neq \bm{0}$ satisfy $\mA \vv = \lambda \vv$. The eigenvalue $\lambda$ tells you how $\mA$ scales along the direction $\vv$.

\subsubsection{Relevance to Deep Learning}

\begin{enumerate}
    \item \textbf{Gradient flow analysis}: Eigenvalues of the Jacobian determine gradient behavior
    \begin{itemize}
        \item $|\lambda| > 1$: Gradients explode
        \item $|\lambda| < 1$: Gradients vanish
        \item $|\lambda| \approx 1$: Stable flow
    \end{itemize}
    
    \item \textbf{Condition number}: $\kappa(\mA) = \frac{\sigma_{max}}{\sigma_{min}}$
    \begin{itemize}
        \item High $\kappa$: Ill-conditioned, unstable optimization
        \item BatchNorm improves conditioning
    \end{itemize}
    
    \item \textbf{PCA}: Eigenvectors of covariance matrix give principal components
\end{enumerate}

\subsection{Singular Value Decomposition (SVD)}

\begin{mathresult}{SVD}
Any matrix $\mA \in \R^{m \times n}$ can be decomposed as:
\begin{equation}
\mA = \mU \bm{\Sigma} \mV^T
\end{equation}
where $\mU \in \R^{m \times m}$ and $\mV \in \R^{n \times n}$ are orthogonal, and $\bm{\Sigma} \in \R^{m \times n}$ is diagonal with singular values $\sigma_1 \geq \sigma_2 \geq \cdots \geq 0$.
\end{mathresult}

\subsubsection{Low-Rank Approximation}

The best rank-$k$ approximation (minimizing Frobenius norm error):
\begin{equation}
\mA_k = \sum_{i=1}^k \sigma_i \vu_i \vv_i^T
\end{equation}

\textbf{Applications}:
\begin{itemize}
    \item LoRA: Low-rank adaptation of weight matrices
    \item Matrix factorization in recommender systems
    \item Compression of embedding matrices
\end{itemize}

\subsection{Norms}

\textbf{Common Norms:}
\begin{align}
\|\vx\|_1 &= \sum_i |x_i| \quad \text{(L1: sparsity-inducing in regularization)} \\
\|\vx\|_2 &= \sqrt{\sum_i x_i^2} \quad \text{(L2: gradient clipping, weight decay)} \\
\|\vx\|_\infty &= \max_i |x_i| \quad \text{(Max norm: constraint-based regularization)}
\end{align}

\textbf{Matrix norms}:
\begin{align}
\|\mA\|_F &= \sqrt{\sum_{i,j} A_{ij}^2} = \sqrt{\tr(\mA^T\mA)} \quad \text{(Frobenius)} \\
\|\mA\|_2 &= \sigma_{max}(\mA) \quad \text{(Spectral)}
\end{align}

%==============================================================================
\section{Calculus for Optimization}
\label{sec:calculus}
%==============================================================================

\subsection{Gradients}

\textbf{Gradient.} For $f: \R^n \to \R$, the gradient $\nabla f(\vx)$ is the vector of partial derivatives. It points in the direction of steepest ascent; gradient descent moves in the opposite direction: $\vx \leftarrow \vx - \eta \nabla f(\vx)$.

\subsection{The Chain Rule}

\begin{mathresult}{Chain Rule}
For composed functions $h = g \circ f$:
\begin{equation}
\frac{dh}{dx} = \frac{dg}{df} \cdot \frac{df}{dx}
\end{equation}

For vector functions $f: \R^n \to \R^m$ and $g: \R^m \to \R^p$:
\begin{equation}
\frac{\partial (g \circ f)}{\partial \vx} = \frac{\partial g}{\partial f} \cdot \frac{\partial f}{\partial \vx}
\end{equation}
(Jacobian multiplication)
\end{mathresult}

\textbf{This is backpropagation!} The backward pass computes:
\begin{equation}
\frac{\partial \mathcal{L}}{\partial \mW_\ell} = \frac{\partial \mathcal{L}}{\partial \vh_L} \cdot \prod_{k=L}^{\ell+1} \frac{\partial \vh_k}{\partial \vh_{k-1}} \cdot \frac{\partial \vh_\ell}{\partial \mW_\ell}
\end{equation}

\subsection{Jacobian and Hessian}

\textbf{Jacobian.} For $f: \R^n \to \R^m$, the Jacobian $\mJ \in \R^{m \times n}$ is the matrix of all partial derivatives. Each row is the gradient of one output. Backpropagation is repeated Jacobian-vector products.

\textbf{Hessian.} For $f: \R^n \to \R$, the Hessian $\mH = \nabla^2 f$ is the matrix of second derivatives. It describes the \textit{curvature} of the loss landscape:
\begin{itemize}
    \item Positive definite $\mH$: Local minimum (bowl-shaped)
    \item Negative definite: Local maximum
    \item Indefinite (mixed eigenvalues): Saddle point
    \item Eigenvalues of $\mH$: Curvature in eigenvector directions
\end{itemize}

\begin{keyinsight}{Why This Matters}
In high-dimensional spaces, most critical points are saddle points, not local minima. SGD naturally escapes saddle points because gradient noise pushes it along negative curvature directions. This is one reason why SGD works well despite non-convex loss landscapes.
\end{keyinsight}

\subsection{Common Derivatives in Deep Learning}

\begin{table}[H]
\centering
\caption{Common derivatives}
\begin{tabular}{ll}
\toprule
\textbf{Function} & \textbf{Derivative} \\
\midrule
$\sigma(x) = \frac{1}{1+e^{-x}}$ & $\sigma(x)(1-\sigma(x))$ \\
$\tanh(x)$ & $1 - \tanh^2(x)$ \\
$\text{ReLU}(x)$ & $\mathbbm{1}[x > 0]$ \\
$\softmax(\vz)_i$ & $\softmax(\vz)_i(\delta_{ij} - \softmax(\vz)_j)$ \\
$\log\softmax(\vz)_i$ & $\delta_{ij} - \softmax(\vz)_j$ \\
$\|\vx\|_2^2$ & $2\vx$ \\
$\vx^T \mA \vx$ & $(\mA + \mA^T)\vx$ \\
\bottomrule
\end{tabular}
\end{table}

%==============================================================================
\section{Probability and Statistics}
\label{sec:probability}
%==============================================================================

\subsection{Probability Fundamentals}

\textbf{Conditional Probability:} $P(A|B) = P(A \cap B) / P(B)$

\begin{mathresult}{Bayes' Theorem}
\begin{equation}
P(\theta | D) = \frac{P(D | \theta) P(\theta)}{P(D)} = \frac{\text{likelihood} \times \text{prior}}{\text{evidence}}
\end{equation}
\end{mathresult}

\textbf{Bayesian interpretation of regularization}:
\begin{itemize}
    \item Prior $P(\theta)$: Regularization term
    \item Gaussian prior → L2 regularization
    \item Laplace prior → L1 regularization
\end{itemize}

\subsection{Expectation and Variance}

\begin{align}
\E[X] &= \sum_x x \cdot P(X=x) = \int x \cdot p(x) dx \\
\Var(X) &= \E[(X - \E[X])^2] = \E[X^2] - \E[X]^2
\end{align}

\textbf{Key identities}:
\begin{align}
\E[aX + b] &= a\E[X] + b \\
\Var(aX + b) &= a^2 \Var(X) \\
\Var(X + Y) &= \Var(X) + \Var(Y) + 2\Cov(X,Y)
\end{align}

\subsection{Important Distributions}

\begin{table}[H]
\centering
\caption{Important distributions in deep learning}
\begin{tabularx}{\textwidth}{lXX}
\toprule
\textbf{Distribution} & \textbf{PDF/PMF} & \textbf{Use in DL} \\
\midrule
Bernoulli & $p^x(1-p)^{1-x}$ & Binary classification \\
Categorical & $\prod_k p_k^{[y=k]}$ & Multi-class classification \\
Gaussian & $\frac{1}{\sqrt{2\pi\sigma^2}}e^{-\frac{(x-\mu)^2}{2\sigma^2}}$ & Regression, VAE \\
Laplace & $\frac{1}{2b}e^{-\frac{|x-\mu|}{b}}$ & L1 prior, robust regression \\
\bottomrule
\end{tabularx}
\end{table}

\subsection{Maximum Likelihood Estimation}

\begin{mathresult}{MLE}
\begin{equation}
\theta_{MLE} = \argmax_\theta P(D|\theta) = \argmax_\theta \sum_{i=1}^N \log P(x_i|\theta)
\end{equation}
\end{mathresult}

\textbf{Connection to cross-entropy}:
For classification with categorical distribution:
\begin{align}
\theta_{MLE} &= \argmax_\theta \sum_i \log P(y_i|x_i; \theta) \\
&= \argmin_\theta -\sum_i \log P(y_i|x_i; \theta) \\
&= \argmin_\theta \mathcal{L}_{CE}
\end{align}

Cross-entropy loss = negative log-likelihood!

%==============================================================================
\section{Information Theory}
\label{sec:information_theory}
%==============================================================================

\subsection{Entropy}

\textbf{Entropy.} Measures average ``surprise'' or uncertainty:
\begin{equation}
H(X) = -\sum_x p(x) \log p(x) = -\E_p[\log p(X)]
\end{equation}

\textbf{Interpretation}: Average ``surprise'' or uncertainty.
\begin{itemize}
    \item $H = 0$: Deterministic (no uncertainty)
    \item $H = \log K$: Uniform over $K$ outcomes (maximum uncertainty)
\end{itemize}

\subsection{Cross-Entropy}

\textbf{Cross-Entropy.} Average surprise when using model $q$ to encode data from true distribution $p$:
\begin{equation}
H(p, q) = -\sum_x p(x) \log q(x) = -\E_p[\log q(X)]
\end{equation}

\textbf{Interpretation}: Average surprise when using $q$ to encode data from $p$.

\textbf{Key property}: $H(p, q) \geq H(p)$, with equality iff $p = q$.

\subsection{KL Divergence}

\begin{mathresult}{KL Divergence}
\begin{equation}
D_{KL}(p \| q) = \sum_x p(x) \log \frac{p(x)}{q(x)} = H(p, q) - H(p)
\end{equation}
\end{mathresult}

\textbf{Properties}:
\begin{itemize}
    \item $D_{KL}(p \| q) \geq 0$ (Gibbs' inequality)
    \item $D_{KL}(p \| q) = 0$ iff $p = q$
    \item \textbf{Not symmetric}: $D_{KL}(p \| q) \neq D_{KL}(q \| p)$
\end{itemize}

\textbf{Applications}:
\begin{itemize}
    \item VAE: KL between posterior and prior
    \item Knowledge distillation: KL between teacher and student
    \item Policy gradient: KL penalty for stability
\end{itemize}

\subsection{Mutual Information}

\textbf{Mutual Information.} How much knowing $Y$ reduces uncertainty about $X$:
\begin{equation}
I(X; Y) = D_{KL}(p(x,y) \| p(x)p(y)) = H(X) - H(X|Y) = H(Y) - H(Y|X)
\end{equation}

\textbf{Interpretation}: How much knowing $Y$ reduces uncertainty about $X$.

\textbf{Properties}:
\begin{itemize}
    \item $I(X; Y) \geq 0$
    \item $I(X; Y) = 0$ iff $X \indep Y$
    \item $I(X; Y) = I(Y; X)$ (symmetric)
\end{itemize}

\textbf{Applications}:
\begin{itemize}
    \item InfoNCE loss: Maximizes lower bound on MI
    \item Information bottleneck: Trade-off compression vs. prediction
    \item Feature selection: High MI with target
\end{itemize}

\subsection{Information Bottleneck Principle}

\begin{keyinsight}[Information Bottleneck]
Good representations should:
\begin{enumerate}
    \item \textbf{Compress}: Minimize $I(X; Z)$ — discard irrelevant information
    \item \textbf{Predict}: Maximize $I(Z; Y)$ — retain task-relevant information
\end{enumerate}

Objective: $\min_Z I(X; Z) - \beta I(Z; Y)$
\end{keyinsight}

This explains the hourglass shape of autoencoders and why intermediate layers have lower dimensionality.

%==============================================================================
\section{Numerical Considerations}
\label{sec:numerical}
%==============================================================================

\subsection{Floating Point Formats}

\begin{table}[H]
\centering
\caption{Floating point formats in deep learning}
\begin{tabular}{lrrrr}
\toprule
\textbf{Format} & \textbf{Bits} & \textbf{Range} & \textbf{Precision} & \textbf{Use} \\
\midrule
FP32 & 32 & $\pm 3.4 \times 10^{38}$ & ~7 digits & Default training \\
FP16 & 16 & $\pm 65504$ & ~3 digits & Mixed precision \\
BF16 & 16 & $\pm 3.4 \times 10^{38}$ & ~2 digits & Modern training \\
INT8 & 8 & $-128$ to $127$ & Exact & Inference \\
\bottomrule
\end{tabular}
\end{table}

\subsection{Numerical Stability Tricks}

\textbf{LogSumExp}:
\begin{equation}
\log \sum_i e^{x_i} = m + \log \sum_i e^{x_i - m}, \quad m = \max_i x_i
\end{equation}

\textbf{Softmax stability}:
\begin{equation}
\softmax(\vx)_i = \frac{e^{x_i - \max(\vx)}}{\sum_j e^{x_j - \max(\vx)}}
\end{equation}

\textbf{Log-space operations}:
\begin{equation}
\log(ab) = \log a + \log b, \quad \log\frac{a}{b} = \log a - \log b
\end{equation}

\begin{warningbox}
\textbf{Common numerical pitfalls}:
\begin{itemize}
    \item $\log(0)$: Always add small $\epsilon$ or use clamp
    \item $\exp(\text{large})$: Use LogSumExp
    \item $0/0$: Check for edge cases in normalization
    \item Gradient of $\sqrt{x}$ at $x=0$: Add $\epsilon$ inside sqrt
\end{itemize}
\end{warningbox}

\begin{warningbox}
\textbf{BF16 vs FP16}: BF16 has the same range as FP32 (8 exponent bits) but lower precision (7 mantissa bits vs 10 in FP16). Use BF16 when you need large dynamic range (gradient accumulation, loss scaling). Use FP16 when you need precision (final inference, metrics computation). Most modern LLM training uses BF16.
\end{warningbox}

%==============================================================================
\section{Interview Questions}
\label{sec:math_interview}
%==============================================================================

\begin{interviewq}
\textbf{Q: Explain SVD and give three applications in machine learning.}

\textbf{What the interviewer is testing:} Whether you understand dimensionality reduction beyond just ``PCA'' and can connect linear algebra to practical ML problems.

\textbf{A:} SVD decomposes any matrix $\mA = \mU \bm{\Sigma} \mV^T$ into orthogonal bases and singular values. Three applications:

\begin{enumerate}
    \item \textbf{LoRA}: Adapts a pretrained weight matrix $\mW$ by adding a low-rank update $\Delta \mW = \mB \mA$ where $\mB \in \R^{d \times r}$ and $\mA \in \R^{r \times d}$ with $r \ll d$. This is motivated by the observation that weight updates during fine-tuning are low-rank.
    \item \textbf{Recommender systems}: Matrix factorization decomposes the user-item interaction matrix into user and item factors---essentially a truncated SVD.
    \item \textbf{Model compression}: Approximate a large weight matrix with a rank-$k$ approximation, reducing parameters from $m \times n$ to $k(m+n)$.
\end{enumerate}

\textbf{Common follow-ups:} ``What's the relationship between SVD and PCA?'' (PCA on centered data = SVD of the data matrix.) ``How does LoRA relate to SVD?'' (LoRA learns the low-rank factors directly rather than computing SVD.)
\end{interviewq}

\begin{interviewq}
\textbf{Q: What is KL divergence? Why is it not a true distance metric? When do you use it in practice?}

\textbf{What the interviewer is testing:} Information theory understanding and practical applications (VAE, knowledge distillation, RLHF).

\textbf{A:} KL divergence $D_{KL}(p \| q)$ measures how much information is lost when using distribution $q$ to approximate $p$. It's not a true distance because it's \textbf{asymmetric}: $D_{KL}(p \| q) \neq D_{KL}(q \| p)$.

The asymmetry matters in practice:
\begin{itemize}
    \item $D_{KL}(p \| q)$ (``forward KL''): Forces $q$ to cover all modes of $p$. Used in VAE (forces posterior to cover prior) and knowledge distillation.
    \item $D_{KL}(q \| p)$ (``reverse KL''): Allows $q$ to mode-seek (ignore some modes of $p$). Used in variational inference, reinforcement learning (PPO's KL penalty).
\end{itemize}

\textbf{Common follow-ups:} ``What's the connection between cross-entropy and KL divergence?'' ($H(p, q) = H(p) + D_{KL}(p \| q)$). ``Why does the VAE use forward KL?''
\end{interviewq}